
intéressant que modulo torsion, i.e. tensorisé par @@M0@@. On voit en même temps, à l'aide de cet exemple, que ce n'est que tensorisé par @@M1@@ que @@M2@@ a un caractère covariant par rapport aux morphismes propres, même en se bornant aux immersions fermées de schémas quasi-projectifs et lisses. En effet, si @@M3@@ est tel que @@M4@@ ne soit pas un isomorphisme, il existe dans @@M5@@ un sous-schéma fermé intègre @@M6@@ tel que @@M7@@ ne soit pas dans @@M8@@, où @@M9@@. Nous pouvons prendre @@M10@@ de dimension minimale, de sorte que @@M11@@. Si @@M12@@ est l'ensemble singulier de @@M13@@, et @@M14@@, @@M15@@, on en conclut que @@M16@@ n'est pas dans @@M17@@ (compte tenu que @@M18@@ est surjectif et compatible avec les augmentations, et que son noyau est @@M19@@). Considérons alors l'inclusion
@@M20@@
qui est une immersion fermée de codimension @@M21@@ de schémas quasi-projectifs et lisses sur @@M22@@, \uline{telle que}
@@M23@@
\uline{ne soit pas compatible avec les} @@M24@@-\uline{filtrations} (\uline{modulo décalage par} @@M25@@), puisque @@M26@@ n'est pas dans @@M27@@.

\textbf{4.7.} Gardant les notations de 4.5, notons maintenant qu'il résulte de (Exp. IX) que le théorème de HODGKIN cité équivaut au fait que l'on a un isomorphisme
@@M28@@
l'homomorphisme @@M29@@ étant à priori surjectif et ayant comme noyau l'idéal engendré par les @@M30@@, où @@M31@@ parcourt les Modules inversibles sur @@M32@@ associés aux représentations @@M33@@ (qui donnent en fait tous les modules inversibles sur @@M34@@ [5, Exp. 15]). On voit d'autre part par un raisonnement analogue que
@@M35@@
est également surjectif, et que son noyau est l'idéal engendré par les éléments de @@M36@@ définis par les représentations @@M37@@, donc en fait par @@M38@@ tout entier.\SGAAuthorFootnote{*}{Cf. IX 3.4, 2.2.} Cela montre que @@M39@@ est réduit à @@M40@@ si et
